\documentclass[a4paper,11pt]{article}
\usepackage[utf8]{inputenc}
\usepackage[T1]{fontenc}
\usepackage[french]{babel}
\usepackage{lmodern}
\usepackage{amsmath,amssymb,amsthm}
\usepackage{mathrsfs}
\usepackage{enumitem}
\usepackage{multicol}

\setlist[enumerate]{label=\textbf{\textcolor{Couleur8}{\arabic*.}}, left=1em}

% Si vous voulez aussi modifier les listes enumerate imbriquées :
\setlist[enumerate,2]{label=\textcolor{Couleur8}{\alph*)}}
\setlist[enumerate,3]{label=\textcolor{Couleur8}{\roman*)}}
\setlist[enumerate,4]{label=\textcolor{Couleur8}{\Alph*)}}
\usepackage{graphicx}
\usepackage{xcolor}
\usepackage{tcolorbox}
\usepackage{geometry}
\usepackage{fancyhdr}
\usepackage{titlesec}
\usepackage{cancel}
\usepackage{ifthen}
\usepackage{tikz}
\usetikzlibrary{arrows.meta}
\usepackage{tkz-tab}
\usepackage{eurosym}

% OPTION PROF/ÉLÈVE
% Décommentez UNE des deux lignes suivantes :
\newboolean{prof}\setboolean{prof}{true}   % Version professeur (avec corrections des devoirs)
\newboolean{prof}\setboolean{prof}{true}  % Version élève (sans corrections des devoirs)

\definecolor{Couleur1}{HTML}{4A5D7A} %Th
\definecolor{Couleur2}{HTML}{A5B5D6} %Prop
\definecolor{Couleur3}{HTML}{A8C68A} %Méthode
\definecolor{Couleur6}{HTML}{8A9CC1} %Def
\definecolor{Couleur7}{HTML}{E6B459} %Rq Voc
\definecolor{Couleur8}{HTML}{D36740} %Titres
\definecolor{correction}{HTML}{D36740} % Couleur pour les corrections
\definecolor{coopmathscorr}{HTML}{F15929} % Couleur corrections coopmaths

% Configuration des marges
\geometry{
	top=2cm,
	bottom=2cm,
	inner=1.5cm,
	outer=1.5cm,
}

% Police
\renewcommand{\familydefault}{\sfdefault}

% Compteurs
\newcounter{seancenum}
\newcounter{seancenumD}
\setcounter{seancenum}{1}
\setcounter{seancenumD}{1}

% Configuration des en-têtes et pieds de page
\pagestyle{fancy}
\fancyhf{}
\ifthenelse{\boolean{prof}}{
    \fancyhead[L]{\textcolor{Couleur8}{\textbf{Corrections Automatismes - Classe de seconde - Version Professeur}}}
}{
    \fancyhead[L]{\textcolor{Couleur8}{\textbf{Corrections Devoirs - Séance 5 - Classe de seconde}}}
}
\fancyhead[R]{\textcolor{Couleur8}{\textbf{2025-2026}}}
\fancyfoot[L]{\textcolor{Couleur8}{Mme Bertail et Mme Pinto}}
\fancyfoot[R]{\textcolor{Couleur8}{\thepage}}
\renewcommand{\headrulewidth}{1pt}
\renewcommand{\headrule}{\hbox to\headwidth{\color{Couleur8}\leaders\hrule height \headrulewidth\hfill}}

% Environnements personnalisés
\newtcolorbox{seance}[1]{
    colback=Couleur6!10,
    colframe=Couleur1!65,
    fonttitle=\bfseries\large,
    title=Correction Séance \theseancenum #1,
    left=5mm,
    right=5mm,
    top=3mm,
    bottom=3mm,
    arc=0mm,
    after=\stepcounter{seancenum}
}

\newtcolorbox{seanceD}[1]{
    colback=Couleur6!5,
    colframe=Couleur6!45,
    fonttitle=\bfseries\large,
    title=Correction Devoirs - Séance \theseancenumD #1,
    left=5mm,
    right=5mm,
    top=3mm,
    bottom=3mm,
    arc=0mm,
    after=\stepcounter{seancenumD}
}

\begin{document}

\setcounter{seancenumD}{5}
\newpage
\begin{seanceD}{} %5

\begin{enumerate}
\item On a : \\
$f(-3)=2\times(-3)^2+(-3)+3=2\times 9-3+3=18-3+3={\color{correction}\boldsymbol{18}}$

Mentalement : On commence par calculer le carré de $-3$, soit $(-3)^2=9$.
On multiplie ensuite cette valeur par le coefficient devant $x^2$, soit $2\times 9=18$.
On ajoute $-3$ soit $18+(-3)=15$.
Pour finir, on ajoute $3$, ce qui donne $15+3$, soit $18$.

\item Comme $g(x)=\dfrac{2}{5}x+6$, on a : \\
$g(-10)=\dfrac{2}{5}\times (-10)+6=2\times \dfrac{-10}{5}+6=-4+6={\color{correction}\boldsymbol{2}}$

\item Comme $u(x)=-2x+1$, on a : \\
$u\left(\dfrac{4}{3}\right)=-2\times \dfrac{4}{3}+1=-\dfrac{8}{3}+1={\color{correction}\boldsymbol{-\dfrac{5}{3}}}$

\item Pour lire l'image de $-3$, on place la valeur de $-3$ sur l'axe des abscisses (axe de lecture des antécédents) et on lit
son image sur l'axe des ordonnées (axe de lecture des images). On obtient : $f(-3)={\color{correction}\boldsymbol{3}}$

\item On a : $g(-1)=-3\times (-1)-6=3-6=-3$ et $g(1)=-3\times 1-6=-3-6=-9$.\\
On en déduit que $g(-1)-g(1)=-3-(-9)={\color{correction}\boldsymbol{6}}$.

\item $\bullet$ Le produit des solutions de l'équation $f(x)=-1$ est égal à $8$.\\  
Cette affirmation est correcte : L'équation $f(x)=-1$ a deux solutions : $-4$ et $-2$.\\
Le produit de ces solutions est $-4\times (-2)= 8$.

La bonne réponse est la réponse {\bfseries \color{correction}B}.
\end{enumerate}

\end{seanceD}

\end{document}